\documentclass[12pt]{article}
\usepackage[utf8]{inputenc}
\usepackage[T1]{fontenc}
\usepackage{hyperref}
\title{Tawbuid Language Resources for K-12 Curriculum Development}
\author{Compiled by Claude Sonnet 5 (claude-sonnet-5) for Andr\'as Kornai}
\date{}
\begin{document}
\maketitle

\section{Introduction}
Tawbuid (ISO 639-3 twb) is a Mangyan language of central Mindoro in the Philippines, with Western (Batangan) and Eastern (Bangon) varieties. Wikipedia gives about 6,800 Western and 1,100 Eastern speakers. The search found no Tawbuid-specific dictionary, textbook or school materials, and the sources disagree on which dialect the Bangon speak. This dossier is short for that reason.

\section{Primary Educational Resources}
None found for Tawbuid.

\section{Linguistic and Reference Materials}
\begin{itemize}
\item Wikipedia, Tawbuid language (overview and dialects). \url{https://en.wikipedia.org/wiki/Tawbuid_language}
\item Omniglot, Tawbuid (Batangan/Bangon): a Latin-alphabet chart and a Batangan prayer text. \url{https://omniglot.com/writing/tawbuid.htm}
\item Wikipedia, Mangyan (context on the Mindoro peoples). \url{https://en.wikipedia.org/wiki/Mangyan}
\end{itemize}

\section{Religious and Community Resources}
The Omniglot page above includes a prayer text. No other religious or community resources were found.

\section{Digital and Audio Resources}
None found for Tawbuid. The UP Diliman Buhid Trilingual Dictionary project is for Buhid, a different Mangyan language, and is not listed as a Tawbuid resource.

\section{Curriculum Development Context}
Tawbuid is a living language, but no documented educational materials turned up. Further searching of SIL, UP Diliman linguistics, and DepEd or NCIP Indigenous Peoples Education materials would be needed before any curriculum work.
\end{document}